\documentclass[a4paper,12pt]{article}
\usepackage{color}
\usepackage{graphicx, gensymb}
\usepackage{amsfonts, cancel, amssymb}
\usepackage{amsmath, amsthm, array, esvect, esint}
\usepackage{typearea}
\usepackage[a4paper,left=2cm,right=2cm,top=2cm,bottom=3cm,bindingoffset=5mm]{geometry}
\newcommand\numberthis{\addtocounter{equation}{1}\tag{\theequation}}
\newcommand{\bra}{}
\newcommand{\kom}{}
\newcommand{\brak}{}
\newcommand{\braket}{}
\newcommand{\intx}{}
\newcommand{\intr}{}
\newcommand{\kla}{}
\newcommand{\klam}{}
\newcommand{\klammer}{}
\newcommand{\oper}{}
\newcommand{\tecxt}{}
\newcommand{\Bra}{}
\newcommand{\Ket}{}

\begin{document}

\title{Helium}
\author{Joachim Schwardt}
\date{\today}
\maketitle

\section{Helium}
\subsection*{a)}
Betrachte in nullter Näherung $H=H_1+H_2$ mit $H_i=-\frac{\hbar^2}{2m}\nabla_i^2-\frac{Ze}{4\pi\epsilon_0r_1}$. Dann lautet die stationäre Schrödinger-Gleichung für Helium:
\begin{align*}
E\phi&=H\phi=H_1\phi +H_2\phi
\end{align*}
\subsection*{b)}
Betrachte den Ansatz $\phi=\phi_1(r_1)\phi_2(r_2)$. Dann folgt für die Energie des Grundzustandes:
\begin{align*}
E\phi&=\phi_2H_1\phi_1+\phi_1H_2\phi_2=\kla\left( {E_1+E_2} \right)\phi\ \ \Rightarrow\ \ E=E_1+E_2\\
&\Rightarrow\ \ E_\tecxt\text{GZ}=-2\cdot\frac{\alpha^2m_ec^2}{2}=-m_ec^2\cdot\frac{e^2}{4\pi\epsilon_0\hbar c}=\underline{-27,21\,\tecxt\text{eV}}
\end{align*}
\subsection*{c)}
Die Korrektur hat einen positiven Wert, da der Wechselwirkungsterm $\frac{e^2}{4\pi\epsilon_0 r_{12}}$ positiv ist. Mit einem Zahlenwert von $E_\tecxt\text{korr.}=3,4\,$eV erhält man dann $E_\tecxt\text{GZ}\approx -23,8\,$eV.
\subsection*{d)}
Die Umrechnung von $\frac{\tecxt\text{J}}{\tecxt\text{mol}}$ in eV erfolgt mit $1\,\tecxt\text{eV}=e\cdot N_A\,\frac{\tecxt\text{J}}{\tecxt\text{mol}}$. Die Ionisierungsenergien sind also $E_1=24.58\,$eV und $E_2=54.41\,$eV, was etwa dem 2-fachen bzw. 4-fachen der Grundzustandsenergie entspricht.
\section{Para- und Orthohelium}
\subsection*{a)}
\subsection*{b)}
Bei Parahelium ist der Grundzustand $1^1\tecxt\text{S}_0$, also $n=1$. Bei Orthohelium $2^3\tecxt\text{S}_1$, also $n=2$. Die Größenordnung der Wellenlänge, die man beim Übergang in diese Grundzustände erwartet ist also:
\begin{align*}
\lambda_\tecxt\text{Para}&=\frac{hc}{E_\tecxt\text{Para}}=\frac{hc\cdot 1^2}{E_R}=\underline{91,1\,\tecxt\text{nm}}\\
\lambda_\tecxt\text{Ortho}&=\frac{hc}{E_\tecxt\text{Ortho}}=\frac{hc\cdot 2^2}{E_R}=\underline{364,5\,\tecxt\text{nm}}
\end{align*}
\section{LS- und jj-Kopplung}
Betrachte ein Zweielektronensystem mit einem $2p-$ und einem $3d-$Elektron. \\
Bei der LS-Kopplung betrachten wir zunächst $L$ und $S$ getrennt. Dann gilt $L=l_1+l_2=1+2=3$ und $S=s_1+s_2=\pm\frac{1}{2}\pm\frac{1}{2}=-1, 0, 0, 1$. Damit folgt $J=2, 3, 3, 4$, es gibt also 4 Zustände.\\
Bei der jj-Kopplung betrachten wir zunächst $j_i$ und bilden dann direkt $J$. Also gilt $j_1=l_1+s_1=1\pm\frac{1}{2}=\frac{1}{2},\frac{3}{2}$ und $j_2=l_2+s_2=2\pm\frac{1}{2}=\frac{3}{2},\frac{5}{2}$. Damit ergibt sich $J=j_1+j_2=2,3,4,3$, also wieder die gleichen 4 Zustände (die Reihenfolge ist hier willkürlich).\\
Betrachte nun zwei Elektronen in der 3p-Schale:
\begin{table}[h!]
\centering
\begin{tabular}{c|c|c|c|c}
$l_1$&$l_2$&$s_1$&$s_2$&Zustand \\ \hline
1&1&+&-&${}^{1}\tecxt\text{D}_{2}$ \\ \hline
 &&+&+&${}^{3}\tecxt\text{P}_{2}$ \\ 
1&0 &+&-&${}^{1}\tecxt\text{P}_{1}$ \\ 
 & &-&-&${}^{-1}\tecxt\text{P}_{0}$ \\ \hline
 & &+&+&${}^{3}\tecxt\text{S}_{1}$ \\ 
1&-1&+&-&${}^{1}\tecxt\text{S}_{0}$ \\ 
 & &-&-&${}^{-1}\tecxt\text{S}_{-1}$ \\ \hline
0&0&+&-&${}^{1}\tecxt\text{S}_{0}$ \\ \hline
 & &+&+&${}^{3}\tecxt\text{P}_{0}$ \\ 
0&-1&+&-&${}^{1}\tecxt\text{P}_{-1}$ \\ 
 & &-&-&${}^{-1}\tecxt\text{P}_{-2}$ \\ \hline
-1&-1&+&-&${}^{1}\tecxt\text{D}_{-2}$
\end{tabular}
\end{table}
\end{document}